% ═══════════════════════════════════════════════════════════════════════
%  Macros et styles TikZ du volume — partagés par le livre (PDF) et
%  les figures autonomes converties en SVG pour le web.
%  Toute macro utilisée dans une formule doit aussi être déclarée dans
%  macros-mathjax.tex (version web).
% ═══════════════════════════════════════════════════════════════════════
% ── Bibliothèques TikZ utilisées dans les figures ────────────────────
\usetikzlibrary{calc, arrows.meta, angles, quotes, patterns, positioning,
                decorations.pathreplacing, decorations.markings, fadings,
                backgrounds, fit}
% Angles, segments et points des figures en coordonnées polaires
\usepackage{tkz-euclide}
\tkzSetUpPoint[fill=black]
\definecolor{fill1}{RGB}{183,183,183}   % remplissage des angles

% ── Raccourcis du volume ─────────────────────────────────────────────
\providecommand{\piup}{\uppi}           % pi droit (anciennes figures)
\newcommand{\boitevide}{\mdlgwhtsquare} % case vide : 5 + □ = 3 (unicode-math)
